\hypersetup{%
	pdftitle = {Modelling single cell dynamics with trajectories and gene regulatory networks},
	pdfsubject = {A subject},
	pdfkeywords = {import keywords},
	pdfauthor = {Robrecht Cannoodt},
	pdfcreator = {\LaTeX},
}

%%%%%%%%%%%%%%%%%%%%%%%%
%%% DEFINE TITLEPAGE %%%
%%%%%%%%%%%%%%%%%%%%%%%%
\begin{titlepage}
	
	\includegraphics[height=1.5cm]{fig/logos/logo_UGent_EN_RGB_color}
	
	\centering
	
	\vfill
	{\Huge \textsf{Modelling single cell dynamics\\with trajectories and gene regulatory networks}}
	
	\vspace{2.25cm}
	
	{\LARGE \textsf{Robrecht Cannoodt}}
	
	\vspace{2.25cm}
	
	{\textsf{Thesis submitted in fulfilment of the requirements for the degree of \\Doctor in Computer Science, 2019}}
	
	\vfill
	
	{\textsf{Supervisors: \\
			Prof. Dr. Yvan Saeys\\Prof. Dr. Katleen De Preter}}
	
	\vfill
	
	\includegraphics[height=.7cm]{fig/logos/icon_UGent_WE_EN_RGB_color}\hfill
	\includegraphics[height=1cm]{fig/logos/vib_rf_inflammation_research_rgb_pos} \hfill
	\includegraphics[height=.8cm]{fig/logos/dambi_logo} \hfill
	\includegraphics[height=.7cm]{fig/logos/cmgg_myshort} \hfill
	\includegraphics[height=.7cm]{fig/logos/FWO_Logo_KleurCMYK} 
	\vspace{1cm}
	
	{\textsf{Vakgroep Toegepaste Wiskunde, Informatica, en Statistiek\\
			Faculteit Wetenschappen, Universiteit Gent \\
			Krijgslaan 281 - S2, 9000 Gent}}
		
\end{titlepage}

%%%%%%%%%%%%%%%%%%%%%%%%
%%% ACKNOWLEDGEMENTS %%%
%%%%%%%%%%%%%%%%%%%%%%%%

\newpage{\thispagestyle{empty}\cleardoublepage}

%%%%%%%%%%%%%
%%% LISTS %%%
%%%%%%%%%%%%%
\newpage{\thispagestyle{empty}\cleardoublepage}
\tableofcontents 


\newpage{\thispagestyle{empty}\cleardoublepage}


%%%%%%%%%%%%%%%%%%%%
%%% NOMENCLATURE %%%
%%%%%%%%%%%%%%%%%%%%
\makenomenclature
\printnomenclature

\nomenclature{HCA}{Human Cell Atlas}
\nomenclature{TCGA}{The Cancer Genome Atlas}

\nomenclature{DNA}{Deoxyribonucleic Acid}
\nomenclature{RNA}{Ribonucleic Acid}
\nomenclature{mRNA}{Messenger RNA}
\nomenclature{TF}{Transcription Factor}

\nomenclature{GRN}{Gene Regulatory Network}
\nomenclature{NI}{Network Inference}
\nomenclature{CART}{Classification And Regression Trees}
\nomenclature{RF}{Random Forests}
\nomenclature{TI}{Trajectory Inference}

\nomenclature{$k$NN}{$k$-nearest-neighbour}
\nomenclature{DR}{Dimensionality Reduction}
\nomenclature{PCA}{Principal Component Analysis}
\nomenclature{MDS}{Multi-Dimensional Scaling}
\nomenclature{t-SNE}{t-distributed Stochastic Neighbor Embedding}
\nomenclature{UMAP}{Uniform Manifold Approximation and Projection}
\nomenclature{LDMS}{Landmark Multi-Dimensional Scaling}
\nomenclature{DE}{Differential Expression}
\nomenclature{TDE}{Trajectory Differential Expression}


\chapter*{Summary}
\vspace{-2em}

Recent developments in the \textit{single-cell omics} technologies triggered a major milestone in the field of (cell)biology. Current techniques can be used to profile both the genome and the transcriptome of individual cells, leading to important insights into cellular heterogeneity and dynamic cell processes. This evolution has led, among other things, to the Human Cell Atlas initiative, an international collaboration project with the aim of mapping all human cell types in terms of its biomolecular contents.

During this doctoral project I have been able to follow these inspiring advancements first hand, but I also made scientific contributions myself, by extensively evaluating new computational tools and by developing various tools and algorithms.

A common problem of pioneering computational tools is the lack of sufficient available datasets to enable the assessment of performances in a quantitative manner. During this PhD, I developed a flexible and extendible \textit{in silico} simulator for individual cells, with the aim of supporting current and future computational research fields and to help kick-start emerging computational sub-domains (Chapter~\ref{chap:dyngen}).

I conducted an extensive study of 45 \textit{trajectory inference} methods (Chapter~\ref{chap:dynbenchmark}). Trajectory inference is a new type of computational analysis with the aim of identifying and studying transitions between main cellular states. While this class of methods is one of the fastest growing subdomains within the \textit{single-cell omics}, a quantitative study of the advantages and disadvantages was hitherto lacking.

In addition, I developed four new software tools. These methods include an extensive toolkit for inferring, visualising and interpreting trajectories (Chapter~\ref{chap:dyno}), a robust tool for inferring linear trajectories (Chapter~\ref{chap:scorpius}), a novel algorithm for studying gene regulation at a single-cell or single-sample level (Chapter~\ref{chap:bred}), and a tool for quantifying the topological changes resulting from minor perturbations in a network (Chapter~\ref{chap:incgraph}).

Finally, I reflect on current challenges within the field, by providing guidelines for how to self-assess trajectory inference methods (Chapter~\ref{chap:selfassessment}) and, more generally, how to perform large-scale benchmarking experiments of computational methods (Chapter~\ref{chap:guidelines}).

Throughout this project I particularly enjoyed working in an open-source ecosystem, collaborating with scientists around the world to tackle unsolved problems, whilst learning best practices by inspecting the source code of exemplary software projects. 


\chapter*{Samenvatting}
\vspace{-8em}

Recente ontwikkelingen in de \textit{single-cell omics} technologieën vormen een belangrijke mijlpaal in het domein van de (cel)biologie. Met de huidige technieken kunnen zowel het genoom als het transcriptoom van individuele cellen geprofileerd worden wat leidt tot belangrijke inzichten in de cellulaire heterogeniteit en de dynamische celprocessen. Deze evolutie heeft onder meer geleid tot het Human Cell Atlas initiatief, een internationaal samenwerkingsproject met als doel het in kaart brengen van alle humane cell types.

Gedurende dit doctoraatsproject heb ik niet alleen deze inspirerende ontwikkelingen van nabij kunnen volgen, maar evenzeer zelf wetenschappelijke bijdragen kunnen aanleveren door het uitgebreid evalueren van nieuwe computationele tools en door het ontwikkelen van verscheiden tools en algoritmen.

Een veelvoorkomend probleem bij de ontwikkeling van innovatieve computationele algoritmen is het gebrek aan voldoende beschikbare datasets om de performantie op een kwantitatieve wijze te kunnen beoordelen. Tijdens dit doctoraat ontwikkelde ik een flexibele en uitbreidbare \textit{in silico} simulator voor individuele cellen, met als doel huidige en toekomstige computationele onderzoeksvelden te ondersteunen en kick-starter projecten een duwtje in de rug te geven (Chapter~\ref{chap:dyngen}). 

Eveneens voerden we een uitgebreide studie van 45 \textit{trajectvoorspellende} methoden uit (Hoofdstuk~\ref{chap:dynbenchmark}). Trajectvoorspellende methoden is een nieuwe soort van computationele analyse met als doel transities tussen verscheidene cellulaire toestanden te identificeren en te bestuderen. Terwijl deze klasse van methoden een van de snelst groeiende subdomeinen binnen de \textit{single-cell omics} vormt, was er tot op dat moment nog geen kwantitatieve studie van de voor- en nadelen van dergelijke methoden uitgevoerd.

Ik ontwikkelde ook zelf vier nieuwe softwareprogramma's, waaronder een uitgebreide toolkit om trajecten te voorspellen, te visualiseren en te interpreteren (Hoofdstuk~\ref{chap:dyno}), een robuust programma om lineaire trajecten te voorspellen (Hoofdstuk~\ref{chap:scorpius}), een nieuwe manier om genregulatie op het niveau van individuele cellen of stalen te bestuderen (Hoofdstuk~\ref{chap:bred}), en een programma om topologische veranderingen ten gevolge van kleine veranderingen in een netwerk te kwantificeren (Hoofdstuk~\ref{chap:incgraph}). 

Ik reflecteer op huidige uitdagingen binnen het veld, door richtlijnen aan te bieden aan ontwikkelaars voor het evalueren van hun trajectvoorspellende methoden (Hoofdstuk~\ref{chap:selfassessment}), en meer algemeen over hoe grootschalige benchmarkingsexperimenten van computationele methoden kunnen worden uitgevoerd (Hoofdstuk~\ref{chap:guidelines}).

Gedurende dit project heb ik vooral genoten van het werken in een textit{open-source} omgeving, en het samenwerken met nationale en internationale collega’s uit het veld om onopgeloste problemen aan te pakken.  Door de broncode van excellente softwareprojecten te bestuderen, leerde ik mezelf de aanbevolen werkwijzen aan.